\documentclass[11pt,letterpaper]{article}


\usepackage[breaklinks, pagebackref=true, urlcolor=blue, colorlinks, citecolor=green!50!black, linkcolor=blue,bookmarksopen,bookmarksnumbered,linktocpage=true]{hyperref}

\usepackage[letterpaper, left=1in, right=1in, top=0.9in, bottom=0.9in]{geometry}
\usepackage[utf8]{inputenc}
\usepackage[american]{babel}
\usepackage[normalem]{ulem}
\usepackage{amsmath, amssymb, cases, amsthm}
\usepackage{thmtools}
\usepackage[shortlabels]{enumitem}
\usepackage[framemethod=TikZ]{mdframed}
\usepackage{bbm}
\usepackage{bm}
\usepackage{microtype}
\usepackage{xcolor}
\usepackage{makecell}
\usepackage{mathtools}
\usepackage{float}
\usepackage{modletters}
\usepackage{comment}
\usepackage{multirow}
\usepackage[sort]{cite}
%\usepackage{natbib}
\usepackage[capitalize,noabbrev]{cleveref}
\usepackage{graphics}
\usepackage{xspace}
% pdfsync omitted: editor synchronization only
\usepackage{tocloft}
\setlength\cftbeforesecskip{0.6em}

\usepackage{footnotehyper}
\makesavenoteenv{tabular}

\usepackage{pifont}
\newcommand{\cmark}{\ding{51}}%
\newcommand{\xmark}{\ding{55}}%

\providecommand{\mleftright}{}\let\mleft\left\let\mright\right

\declaretheorem[numberwithin=section,refname={Theorem,Theorems},Refname={Theorem,Theorems}]{theorem}
\declaretheorem[numberwithin=section,refname={Theorem,Theorems},Refname={Theorem,Theorems}]{thm}
\declaretheorem[numberlike=theorem]{lemma}
\declaretheorem[numberlike=theorem]{proposition}
\declaretheorem[numberlike=theorem]{corollary}
\declaretheorem[numberlike=theorem,style=definition]{definition}
\declaretheorem[numberlike=theorem, refname={Claim, Claims}, Refname={Claim, Claims}]{claim}
\declaretheorem[numberlike=theorem,style=remark]{remark}
\declaretheorem[numberlike=theorem,refname={Fact,Facts},Refname={Fact,Facts},name={Fact}]{fact}
\declaretheorem[numberlike=theorem, refname={Question,Questions},Refname={Question,Questions},name={Question}]{question}
\declaretheorem[numberlike=theorem, refname={Problem,Problems},Refname={Problem,Problems},name={Problem}]{problem}
\declaretheorem[numberlike=theorem, refname={Observation,Observations},Refname={Observation,Observations},name={Observation}]{observation}
\declaretheorem[numberlike=theorem, refname={Experiment,Experiments},Refname={Experiment,Experiments},name={Experiment}]{experiment}
\let\primitive\relax
\declaretheorem[numberlike=theorem,refname={Primitive,Primitives},Refname={Primitive,Primitives},name={Primitive}]{primitive}
\declaretheorem[numberlike=theorem,refname={Technique,Techniques},Refname={Technique,Techniques},name={Technique}]{technique}
\declaretheorem[numberlike=theorem,refname={Conjecture,Conjectures},Refname={Conjecture,Conjectures},name={Conjecture}]{conjecture}

\def\cameraready{0}  
\def\final{0}  
\def\iflong{\iffalse}
\ifnum\final=0  

\newcommand{\tawei}[1]{{\color{blue} [{\bf Ta-Wei:} #1]}}
\newcommand{\parth}[1]{{\color{purple} [{\bf Parth:} #1]}}
\newcommand{\TODO}[1]{{\color{blue!50!black} [{\bf Todo:} #1]}}
\else 
\newcommand{\yonggang}[1]{}
\newcommand{\danupon}[1]{}
\newcommand{\sagnik}[1]{}
\newcommand{\todo}[1]{}
\newcommand{\yuval}[1]{}
\newcommand{\jan}[1]{}
\newcommand{\blikstad}[1]{}
\newcommand{\tawei}[1]{}
\newcommand{\parth}[1]{}
\newcommand{\TODO}[1]{}
\fi  

\bibliographystyle{alpha}

\DeclarePairedDelimiter{\ceil}{\lceil}{\rceil}
\newcommand{\eps}{\varepsilon}

\newcommand{\mindeg}{\mathrm{mindeg}}
\newcommand{\nbh}{\mathrm{nbh}}
\newcommand{\spf}{\mathrm{spf}}
\newcommand{\direct}{MDCP }
\newcommand{\best}{\mathsf{best}}
\newcommand{\polylog}{\mathrm{polylog}}
\newcommand{\poly}{\mathrm{poly}}
\newcommand{\outdeg}{\mathrm{outdeg}}
\newcommand{\ones}{\mathrm{ones}}
\DeclareMathOperator*{\var}{Var}
\newcommand{\euler}{\mathrm{e}}
\newcommand{\set}[2][ ]{\{#2 \ifthenelse{\equal{#1}{ }}{ }{~|~#1}\}}
\newcommand{\fvedge}{\mathsf{FindViolatingEdge}}
\newcommand{\XORQ}{\mathsf{XOR}}
\newcommand{\ANDQ}{\mathsf{AND}}

\newcommand{\cc}{\mathrm{cc}}
\newcommand{\Rcal}{\mathcal{R}}
\newcommand{\Pcal}{\mathcal{P}}
\newcommand{\Acal}{\mathcal{A}}
\newcommand{\Scal}{\mathcal{S}}
\newcommand{\Qcal}{\mathcal{Q}}
\newcommand{\pflag}{\mathrm{pflag}}
\newcommand{\R}{\mathbb{R}}
\newcommand{\N}{\mathbb{N}}
\newcommand{\one}{\mathbf{1}}
\newcommand{\inact}{\mathsf{Inactive}}
\newcommand{\act}{\mathsf{Active}}
\newcommand{\cut}{\mathrm{cut}}
\newcommand{\cutd}{\overrightarrow{\mathrm{cut}}}
\newcommand{\Ed}{\overrightarrow{E}}
\newcommand{\mincut}{\lambda}
\newcommand{\tOh}{\widetilde{O}}
\newcommand{\ith}{i^{\scriptsize \mbox{{\rm th}}}}
\newcommand{\jth}{j^{\scriptsize \mbox{{\rm th}}}}



\newcommand{\intpot}[1]{Int_{pot}(#1)}


\newcommand{\ip}[1]{\mleft\langle #1\mright\rangle }
\newcommand{\expt}[2]{\underset{#1}{\mathbb{E}}\mleft[#2\mright]}
\newcommand{\A}{{\mathcal{A}}}
\newcommand{\BPM}{\mathsf{BPM}}
\newcommand{\BMM}{\mathsf{BMM}}
\newcommand{\UBPM}{\mathsf{UBPM}}

\newcommand{\T}{{\mathcal{T}}}

\newcommand{\U}{{\mathcal{U}}}

\newcommand{\bcc}[1]{\textbf{BCAST($#1$)}}
\newcommand{\dw}[1]{#1^{\downarrow}}

\newcommand{\MM}{\mathcal{P}}
\newcommand{\ORQ}{\mathsf{OR}}
\newcommand{\ISQ}{\mathsf{IS}}
\newcommand{\vol}{\mathrm{vol}}
\newcommand{\HPC}{\mathsf{HPC}}
\newcommand{\distance}{\mathsf{dist}}

\newcommand{\rank}{\mathsf{rank}}
\newcommand{\sspan}{\mathsf{span}}
\DeclareMathOperator*{\E}{\mathbb{E}}


\usepackage[most]{tcolorbox}
\tcbuselibrary{skins,breakable}
\tcbset{enhanced jigsaw}
\newcounter{myexample}
\usepackage{xparse}
\usepackage{lipsum}

\makeatletter
\newcommand\footnoteref[1]{\protected@xdef\@thefnmark{\ref{#1}}\@footnotemark}
\makeatother


\renewcommand{\paragraph}[1]{\medskip\noindent{\bf #1}\xspace}






\newcommand{\ground}{U}


\usepackage{tikz}
\usetikzlibrary{arrows}
\usetikzlibrary{arrows.meta}
\usetikzlibrary{shapes}
\usetikzlibrary{backgrounds}
\usetikzlibrary{positioning}
\usetikzlibrary{decorations.markings}
\usetikzlibrary{patterns}
\usetikzlibrary{calc}
\usetikzlibrary{fit}
\usetikzlibrary{decorations}


\renewcommand{\geq}{\geqslant}
\renewcommand{\leq}{\leqslant}
\renewcommand{\le}{\leq}
\renewcommand{\ge}{\geq}
% Parth's macros

\theoremstyle{definition}
\newtheorem{algocf}{Algorithm}
\newenvironment{Algorithm}{\begin{abox}\begin{algocf}}{\end{algocf}\end{abox}}

\Crefname{algocf}{Algorithm}{Algorithms}

\newenvironment{abox}{\begin{tcolorbox}[
		enlarge top by=5pt,
		enlarge bottom by=5pt,
		breakable,
		frame hidden,
		overlay broken = {
			\draw[line width=1pt, black]
			(frame.north west) rectangle (frame.south east);},
		overlay = {
			\draw[line width=1pt, black]
			(frame.north west) rectangle (frame.south east);},
		boxsep=0pt,
		left=4pt,
		right=4pt,
		top=10pt,
		arc=0pt,
		boxrule=1pt,toprule=1pt,
		colback=white
	]}
{\end{tcolorbox}}

\newcommand{\degen}{\textnormal{\textsf{DEGEN}}}

\newcommand{\fail}{\textsc{Reject}}
\newcommand{\true}{\textsc{Accept}}



\providecommand{\given}

\DeclarePairedDelimiter{\pparen}()
\DeclarePairedDelimiter{\card}\vert\vert

\DeclarePairedDelimiterX{\brac}[1][]{
%\renewcommand\given{\nonscript\: \delimsize\vert{} \nonscript\: \mathopen{}}
\ifblank{#1}{\,\cdot\,}{#1}}

\DeclarePairedDelimiterX{\Set}[1]\{\}{
  \renewcommand\given{\nonscript\: \delimsize\vert{} \nonscript\: \mathopen{}}
\ifblank{#1}{\,\cdot\,}{#1}}

\renewcommand{\given}{\nonscript\: \vert \nonscript\:}


\NewDocumentCommand{\Prob}{sO{}E{_}{{}}m}{%
    \Pr_{#3}
    \IfBooleanTF{#1}
    {\brac*{#4}}
    {\brac[#2]{#4}}
}

\DeclarePairedDelimiterXPP{\Ot}[1]{\widetilde{O}}(){}{#1}
\DeclarePairedDelimiterXPP{\Omgt}[1]{\widetilde{\Omega}}(){}{#1}

\DeclarePairedDelimiterXPP{\Var}[1]{\var}[]{}{#1}

\newcommand{\dsi}{\cD_\SI}
\newcommand{\pisi}{\pi_\SI}

\newcommand{\ptd}{\Lambda}
\DeclarePairedDelimiterXPP{\PTD}[2]{\Lambda}(){}{#1, #2}

\DeclareMathOperator*{\distrib}{dist}
\NewDocumentCommand{\Dist}{sO{}E{_}{{}}m}{%
  \mu_{#3}
  \IfBooleanTF{#1}
  {\pparen*{#4}}
  {\pparen[#2]{#4}}
}

\DeclareMathOperator*{\expect}{\mathbb{E}}
\NewDocumentCommand{\Exp}{sO{}E{_}{{}}m}{%
  \expect_{#3}
  \IfBooleanTF{#1}
  {\brac*{#4}}
  {\brac[#2]{#4}}
}

\newcommand{\re}{\rv{e}}

\newcommand{\estar}{e^{\ast}}
\newcommand{\restar}{\re^{\ast}}

\newcommand{\rPi}{\rv{\Pi}}

\newcommand{\Perm}{\Gamma}
\newcommand{\rPerm}{\rv{\Gamma}}

\newcommand{\uni}{\mathcal{U}}

\newcommand{\qpi}{q^{\Pi}}

\newcommand{\score}{\mathsf{Score}}
\newcommand{\totscore}{\mathsf{TScore}}


\DeclareMathOperator*{\inform}{\mathbb{I}}
\DeclarePairedDelimiterXPP{\Inf}[1]{\inform}(){}{#1}

\DeclareMathOperator*{\entropy}{\mathbb{H}}
\DeclarePairedDelimiterXPP{\Ent}[1]{\entropy}(){}{#1}

\DeclareMathOperator{\myIC}{\mathsf{IC}}
\DeclareMathOperator{\myCC}{\mathsf{CC}}
\newcommand{\ExtIC}{\myIC^{\textnormal{ext}}}
\newcommand{\IntIC}{\myIC^{\textnormal{int}}}

\newcommand{\rPTD}{\mathbf{\Lambda}}

\newcommand{\Leq}[1]{{\underset{\textnormal{#1}}\leq}}
\newcommand{\Geq}[1]{{\underset{\textnormal{#1}}\geq}}
\newcommand{\Eq}[1]{{\underset{\textnormal{#1}}=}}



\newcommand{\bhpc}{\ensuremath{\textnormal{\textsf{BHPC}}}\xspace}
\newcommand{\bmhpc}{\ensuremath{\textnormal{\textsf{BMHPC}}}\xspace}
\newcommand{\hpc}{\ensuremath{\textnormal{\textsf{HPC}}}\xspace}
\newcommand{\mhpc}{\ensuremath{\textnormal{\textsf{MHPC}}}\xspace}
\newcommand{\setint}{\ensuremath{\textnormal{\textsf{SetInt}}}\xspace}
\newcommand{\defc}{\mathtt{def}}
\newcommand{\odeg}{\text{odeg}}
\newcommand{\Vaux}{V_{\text{aux}}}
\newcommand{\Eaux}{E_{\text{aux}}}
\newcommand{\prantar}[1]{{\color{blue} [{\bf Prantar:} #1]}}
\newcommand{\mypar}[1]{\smallskip\noindent{\sffamily\bfseries #1.}~}
\newcommand{\floor}[1]{\lfloor #1 \rfloor}




\newcommand{\cmax}{\ensuremath{\textnormal{\textsf{c}}_{\textnormal{rs}}}}
\newcommand{\protMO}{\ensuremath{\Prot_{\mo}}\xspace}
\newcommand{\bProtMO}{\ensuremath{\bm{\Prot}_{\mo}}\xspace}
\newcommand{\protBHH}{\ensuremath{\Prot_{\bhh}}\xspace}
\newcommand{\bProtBHH}{\ensuremath{\bm{\Prot}_{\bhh}}\xspace}


\newcommand{\Paren}[1]{\Big(#1\Big)}
\newcommand{\Bracket}[1]{\Big[#1\Big]}
\newcommand{\bracket}[1]{\mleft[#1\mright]}
\newcommand{\paren}[1]{\ensuremath{\mleft(#1\mright)}\xspace}

\newenvironment{tbox}{\begin{tcolorbox}[
		enlarge top by=5pt,
		enlarge bottom by=5pt,
		 breakable,
		 boxsep=0pt,
                  left=4pt,
                  right=4pt,
                  top=10pt,
                  boxrule=1pt,toprule=1pt,
                  colback=white,
                  arc=-1pt,
        %          drop shadow={black,opacity=1},
                 % arc=0pt
                  ]%%
	}
{\end{tcolorbox}}



\newcommand{\Prot}{\ensuremath{\Pi}}
\newcommand{\prot}{\ensuremath{\pi}}
\newcommand{\Rpub}{{R}^{pub}}
\newcommand{\Rpri}{{R}^{pri}}

\newcommand{\bProt}{\bm{\Prot}}
\newcommand{\bRpub}{\bm{\Rpub}}
\newcommand{\bRpri}{\bm{\Rpri}}


\newcommand{\bV}{\bm{V}}
\newcommand{\bX}{\bm{X}}
\newcommand{\bY}{\bm{Y}}
\newcommand{\bZ}{\bm{Z}}
\newcommand{\bA}{\bm{A}}
\newcommand{\bP}{\bm{P}}
\newcommand{\bR}{\bm{R}}
\newcommand{\bM}{\bm{M}}
\newcommand{\bI}{\bm{I}}
\newcommand{\bJ}{\bm{J}}
\newcommand{\bSigma}{\bm{\sigma}}
\newcommand{\bG}{\bm{G}}

\newcommand{\bB}{\ensuremath{\bm{B}}}
\newcommand{\bC}{\ensuremath{\bm{C}}}
\newcommand{\bD}{\ensuremath{\bm{D}}}

\newcommand{\supp}[1]{\ensuremath{\textsc{supp}(#1)}}




\newtheorem{mdresult}{Result}
\newenvironment{result}{\begin{mdframed}[backgroundcolor=lightgray!40,topline=false,rightline=false,leftline=false,bottomline=false,innertopmargin=2pt]\begin{mdresult}}{\end{mdresult}\end{mdframed}}


%-------------------------- Textbox ---------------------------------

\newcommand{\textbox}[2]{
{
\begin{tbox}
\textbf{#1}
{#2}
\end{tbox}
}
}

\newcommand{\rA}{\rv{A}}
\newcommand{\rB}{\rv{B}}
\newcommand{\rC}{\rv{C}}
\newcommand{\rD}{\rv{D}}

\newcommand{\rX}{\rv{X}}
\newcommand{\rY}{\rv{Y}}
\newcommand{\rI}{\rv{I}}

\newcommand{\rv}[1]{\ensuremath{\mathsf{#1}}}

\newcommand{\II}{\ensuremath{\mathbb{I}}}
\newcommand{\HH}{\ensuremath{\mathbb{H}}}
\newcommand{\mi}[2]{\ensuremath{\II(#1 \,; #2)}}
\newcommand{\en}[1]{\ensuremath{\HH(#1)}}

\newcommand{\itfacts}[1]{Fact~\ref{fact:it-facts}-(\ref{part:#1})\xspace}

\newcommand{\DD}{\ensuremath{\mathbb{D}}}

\newcommand{\distribution}[1]{\ensuremath{\textnormal{dist}(#1)}\xspace}

\newcommand{\XX}{\ensuremath{\mathcal{X}}}
\newcommand{\YY}{\ensuremath{\mathcal{Y}}}


\newcommand{\HPCk}{\ensuremath{\HPC_k}\xspace}

\newcommand{\protHPC}{\ensuremath{\prot_{\HPC}}}

\newcommand{\distHPC}{\ensuremath{\dist_{\HPC}}}

\newcommand{\unif}{\ensuremath{\mathcal{U}}}

\newcommand{\rZ}{\rv{Z}}
\newcommand{\rt}{\rv{t}}
\newcommand{\rT}{\rv{T}}
\newcommand{\rE}{\rv{E}}
\newcommand{\rP}{\rv{P}}
\newcommand{\rProt}{\rv{\Prot}}

\newcommand{\distSI}{\ensuremath{\dist}_{\ensuremath{\textnormal{\textsf{SI}}}}\xspace}


\newcommand{\protSI}{\ensuremath{\prot}_{\ensuremath{\textnormal{\textsf{SI}}}}\xspace}
\newcommand{\ProtSI}{\ensuremath{\Prot}_{\ensuremath{\textnormal{\textsf{SI}}}}\xspace}
\newcommand{\rProtSI}{\ensuremath{\rv{\Prot}}_{\ensuremath{\textnormal{\textsf{SI}}}}\xspace}

\newcommand{\distPI}{\ensuremath{\dist}_{\ensuremath{\textnormal{\textsf{PI}}}}\xspace}
\newcommand{\protPI}{\ensuremath{\prot}_{\ensuremath{\textnormal{\textsf{PI}}}}\xspace}
\newcommand{\ProtPI}{\ensuremath{\Prot}_{\ensuremath{\textnormal{\textsf{PI}}}}\xspace}
\newcommand{\rProtPI}{\ensuremath{\rv{\Prot}}_{\ensuremath{\textnormal{\textsf{PI}}}}\xspace}


\newcommand{\SI}{\ensuremath{\textnormal{\textsf{Set-Int}}}\xspace}
\newcommand{\PI}{\ensuremath{\textnormal{\textsf{Pair-Int}}}\xspace}


\newcommand{\errs}{\ensuremath{\textnormal{errs}}\xspace}




\newcommand{\rK}{\rv{K}}
\newcommand{\rJ}{\rv{J}}
\newcommand{\rL}{\rv{L}}
\newcommand{\rS}{\rv{S}}
\newcommand{\rM}{\rv{M}}
\newcommand{\rR}{\rv{R}}
\newcommand{\rbarS}{\rv{\barS}}



\newcommand{\distDisjN}{\ensuremath{\distDisj^{\textnormal{\textsf{N}}}}\xspace}
\newcommand{\distDisjY}{\ensuremath{\distDisj^{\textnormal{\textsf{Y}}}}\xspace}

\newcommand{\mii}[3]{\ensuremath{\II_{#3}(#1 \,; #2)}}

\newcommand{\disjin}[4]{ \ensuremath{\begin{bmatrix} #1&\hspace{-8pt}#2 \\ #3&\hspace{-10pt}#4 \end{bmatrix}}}
\renewcommand{\disjin}[4]{\textnormal{[\ensuremath{#1#2,#3#4}]}}

\newcommand{\barS}{\overline{S}}

\newcommand{\hd}[2]{\ensuremath{\textnormal{h}(#1,#2)}\xspace}
\newcommand{\hdt}[2]{\ensuremath{\textnormal{h}^2(#1,#2)}\xspace}


\newcommand{\CC}[2]{\ensuremath{\textnormal{\textsf{CC}}_{#2}(#1)}\xspace}
\newcommand{\ICost}[2]{\ensuremath{\textnormal{\textsf{IC}}_{#2}(#1)}\xspace}
\newcommand{\IC}[2]{\ICost{#1}{#2}}


\newcommand{\rTheta}{\rv{\Theta}}



\newcommand{\comp}{\succ_{\ProtSI}}
\newcommand{\compr}{\prec_{\ProtSI}}
\newcommand{\factor}{3}
%%%%
\title{Near-linear deterministic communication protocol for exact degeneracy and k-cores}
\author{Sepehr Assadi, Prantar Ghosh, Bruno Loff, Parth Mittal, Sagnik Mukhopadhyay}
\date{doi:10.4230/LIPIcs.CCC.2024.7; CC BY 4.0}
\begin{document}
\maketitle
\setcounter{section}{6}
\noindent\textit{Editorial note.} Native author TeX from fixed arXiv v1, distributed under CC BY 4.0. All proof-bearing communication protocols retained, including the warm-up as supplied by the authors. Source-local graph definitions and standard peeling facts are included. Front matter and unrelated lower bounds omitted. Counted result is thm:communication-upper only.\par\medskip
\noindent\textit{Allowed standard prerequisites (editorial).}\begin{itemize}
\item Classical peeling characterization of degeneracy. Matula and Beck 1983, author source fact:min-k-ord and fact:peelingmin-degen retained
\end{itemize}\medskip
\subsection{\texorpdfstring{$k$-Cores}{k-Cores} and Degeneracy in Graph Streams}\label{sec:degen-intro}

For any undirected graph $G=(V,E)$ and integer $k \geq 1$, a \emph{$k$-core} in $G$ is a maximal set $S$ of vertices such that the induced subgraph of $G$ on $S$, denoted by $G[S]$, has a minimum degree of at least $k$. In other words, 
any vertex in $S$ has at least $k$ other neighbors in $S$. 

$k$-Cores provide a natural notion of well-connectedness in massive graphs, and as such, computing $k$-cores (and more generally $k$-core decompositions; see, e.g.,~\cite{LiuSYDS22}) 
has been widely studied in databases~\cite{BonchiGKV14, ChuZ00ZXZ20, LiZZQZL19}, social networks~\cite{DhulipalaBS17, DhulipalaBS18, KhaouidBST15}, machine learning~\cite{Alvarez-HamelinDBV05, EsfandiariLM18, GhaffariLM19}, 
among others~\cite{LiWSX18,GalimbertiBGL20, SariyuceGJWC13}. 

As a result, in recent years, there has been a rapidly growing body of work on
computing $k$-cores on massive graphs in parallel and streaming
models of computation~\cite{SariyuceGJWC13,DhulipalaBS17, DhulipalaBS18,EsfandiariLM18, GhaffariLM19,LiuSYDS22}. In particular,~\cite{EsfandiariLM18} presented a single-pass algorithm that for any $\eps > 0$, 
computes a $(1-\eps)$-approximation of every $k$-core in $G$ (i.e., obtains a $(1-\eps)$-approximate $k$-core decomposition) in $\Ot{n/\eps^2}$ space (see also~\cite{SariyuceGJWC13} for an earlier streaming algorithm and~\cite{GhaffariLM19}
for a closely related parallel algorithm). 

The \emph{degeneracy} of a graph $G=(V,E)$, denoted by $\kappa(G)$, is the largest integer $k \geq 0$ such that $G$ contains a non-empty $k$-core. 
The simple greedy algorithm that at every step peels off the smallest degree vertex results in the so-called \emph{degeneracy ordering} of $G$ and $\kappa{(G)}$ 
is equal to the largest degree of a vertex removed in this peeling process~\cite{MatulaB83}. Degeneracy is a standard measure of uniform sparsity and is closely
related to other such notions like arboricity (which is within a factor $2$ of degeneracy). Moreover, computing degeneracy is a subroutine for approximating various
other problems such as arboricity~\cite{AlonYZ97}, densest subgraph~\cite{Charikar00}, $(\kappa+1)$ coloring~\cite{erdHos1966chromatic}. 

The degeneracy problem, and closely related uniform-sparsity measures such as densest subgraph, have also been studied extensively in the graph 
streaming literature~\cite{BahmaniKV12,Farach-ColtonT14,BhattacharyaHNT15,McGregorTVV15,Farach-ColtonT16,BeraCG19,AlonA20}. In particular,~\cite{Farach-ColtonT14} provided an $O(\log{n})$-pass semi-streaming algorithm that outputs a constant factor approximation to degeneracy, and~\cite{Farach-ColtonT16} subsequently improved this to a single-pass $(1-\eps)$-approximation in $\Ot{n/\eps^2}$ space (see also~\cite{McGregorTVV15} for densest subgraph and~\cite{BeraCG19,AlonA20} for degeneracy coloring). 
 
\medskip

In terms of lower bounds, \cite{BeraCG19} prove that any \emph{single-pass} streaming algorithm that computes the exact value of degeneracy or approximates it within an additive factor of $\lambda$ requires $\Omega(n^2)$ space or $\Omega(n^2/\lambda^2)$ space respectively. Our \emph{polynomial-pass} lower bounds for $k$-cores and degeneracy, now in a very strong sense, rule out the possibility of extending any prior semi-streaming algorithms computing near-optimal solutions to these problems, to compute \emph{exactly} optimal solutions. 


\subsection{Degeneracy and \texorpdfstring{$k$-Cores}{k-Cores}}

The definition of a $k$-core and the degeneracy $\kappa(G)$ of a graph $G$ can be found in \Cref{sec:degen-intro}. Here, we define some related terms that come up often in the paper. We also state some basic facts about them that we later use in our proofs. 


Given an ordering $\sigma$ of vertices, we say $v\prec_\sigma u$ if $u$ appears (somewhere) after $v$ in $\sigma$. 

\begin{definition}[Outdegree with respect to an ordering]\label{def:odegordering}
    For a graph $G=(V,E)$. given an ordering $\sigma$ of vertices in $V$, the outdegree of a vertex $v$ with respect to $\sigma$, denoted by $\odeg_{\sigma}(v)$, is the number of neighbors of $v$ that appear after it in $\sigma$. Formally,
$$\odeg_{\sigma}(v) = |\{u\in N(v): v \prec_\sigma u\}|$$ 
\end{definition}

We can also think of it as orienting all edges $(x,y)$ from $x$ to $y$ if $x\prec_\sigma y$. Then $\odeg_\sigma(v)$ is the outdegree of $v$ in the resultant oriented graph. 


\begin{definition}[$k$-ordering]\label{def:k-ordering} 
    Given a graph $G=(V,E)$, an ordering $\sigma$ of vertices in $V$ is called a $k$-ordering if $\forall v\in V: \odeg_\sigma(v)\leq k$.   
\end{definition}

The following fact connects vertex orderings to degeneracy. 

\begin{fact}[\cite{MatulaB83}]\label{fact:min-k-ord}
    The minimum number $k$ for which a graph $G$ admits a $k$-ordering is $\kappa(G)$.
\end{fact}


The above fact implies that to prove that a graph has degeneracy at most $k$, it suffices to exhibit a $k$-ordering of the nodes. We now formally define \emph{degeneracy ordering} of a graph.  

\begin{definition}[Degeneracy ordering] \label{def:degen-ord}
    Given a graph $G=(V,E)$, a $\kappa(G)$-ordering of $V$ is called a degeneracy ordering of $G$.
\end{definition}

A degeneracy ordering of a graph can be obtained by running the following ``peeling algorithm'' on it: starting with an empty sequence $\sigma$, recursively remove the minimum degree vertex from $G$ and append it to $\sigma$ until the graph becomes empty. Again, the maximum degree that a node has when it is peeled by this algorithm is exactly $\kappa(G)$. If not, then the peeling algorithm gives a $k$-ordering for some $k<\kappa(G)$ which violates \Cref{fact:min-k-ord}. Hence, we get the following fact. 


 
  

\begin{fact}\label{fact:peelingmin-degen}
Suppose the peeling algorithm on a graph $G$ removes its nodes in the order $v_1,\ldots,v_n$. Define $G_0:=G$ and $G_i:= G\setminus \{v_1,\ldots,v_i\}$ for $i\in [n]$. Then, $$\kappa(G)=\max_{i\in [n]} \deg_{G_{i-1}}(v_i)$$. 
\end{fact}




\section{Communication Upper Bounds}\label{sec:communication-upper}

In this section, we will see a communication upper bound for computing
the degeneracy of a graph exactly.
The setting is as follows:
Alice and Bob are respectively given \emph{disjoint} edge sets $E_A$ and $E_B$ on the set of
vertices $V$ (known to both Alice and Bob), and wish to compute the degeneracy of the graph
$G = (V, E_A \cup E_B)$.
The main theorem of this section is:


\begin{theorem}\label{thm:communication-upper}
  There exists a deterministic communication protocol that, given an $n$-vertex
  graph $G$, computes the degeneracy $\kappa(G)$ using $O(n \log^3 n)$ bits of
  communication.
  Moreover, the protocol outputs a $\kappa(G)$-ordering of the vertices
  of $G$, and a $\kappa(G)$-core in $G$.
\end{theorem}

We will focus on a decision version of the degeneracy problem, which we define
in the lemma below:

\begin{lemma}\label{lem:degen-decision}
  There exists a deterministic communication protocol that given an $n$-vertex
  graph $G$ and an integer $k$, decides whether the degeneracy
  $\kappa(G) \leq k$ using $O(n \log^2 n)$ bits of communication.
  Moreover, if the protocol accepts, it outputs a $k$-ordering of the vertices
  of $G$, and if it rejects, it outputs a $(k + 1)$-core in $G$.
\end{lemma}

For brevity, we will call the decision version of degeneracy defined in the
lemma above $\degen_{n, k}$.
Note that once we have a protocol for the decision problem, a protocol for the
search variant follows via binary search.
\begin{proof}[Proof of \Cref{thm:communication-upper}]
  Alice and Bob can binary search for the minimum $k$ such that $G$ has
  degeneracy at most $k$ (which is exactly $\kappa(G)$) using the protocol
  given by \Cref{lem:degen-decision}.
  Note that this incurs a multiplicative $\log n$ overhead in communication
  cost.
  Further, the protocol in the lemma also gives a $\kappa(G)$-ordering, since
  it accepts the input $(G, \kappa(G))$.
  Finally, to obtain a $\kappa(G)$-core, Alice and Bob can run the protocol
  in the lemma on the input $(G, \kappa(G) - 1)$.
\end{proof}

In the rest of this section, we will prove \Cref{lem:degen-decision}.
First, we show an easier (but more expensive) protocol that captures the main
idea.

\subsection{Warm-Up: An \texorpdfstring{$\Ot{n\sqrt{n}}$}{Õ(n sqrt n)} Communication Protocol}\label{sec:nrootn}
Recall the classical peeling algorithm, where one successively peels off the
vertex of minimum degree to find a degeneracy ordering of $G$.
We state a decision version of it below, which removes an arbitrary vertex of
degree $\leq k$ in each step.

\begin{Algorithm}\label{alg:peeling}
  The Peeling Algorithm. \medskip

  \textbf{Input:} A graph $G$, and an integer $k$.\\
  \textbf{Output:} \true{} if $G$ has degeneracy $\leq k$, \fail{} otherwise.

  \begin{enumerate}
    \item While $G$ is non-empty:
      \begin{enumerate}
        \item If there is a vertex $v \in G$ of degree $\leq k$,
          remove $v$ from $G$.
        \item Else output \fail.
      \end{enumerate}
    \item Output \true.
  \end{enumerate}
  %  \caption{The Peeling Algorithm for $\degen(G, k)$}\label{alg:peeling}
\end{Algorithm}

Note that in the \true{} case the $k$-ordering is simply the order the vertices
were peeled in, and in the \fail{} case, the $(k + 1)$-core is the set of vertices
remaining in $G$ when there are no vertices of degree $\leq k$.

We want to somehow port this algorithm to the communication setting.
The immediate problem is that Alice or Bob cannot compute a low degree
vertex without any communication after a vertex is removed.
The naive way to get around this is to send the updated degree of each
vertex after each deletion, but this can cost $\Omega(n^2)$ communication.
Instead, we observe that we only need to update the degrees of a remaining
vertex if it is not too large:

\begin{observation}\label{obs:sqrt-local}
  If $\deg(u) > k + \sqrt{n}$ at some step of the peeling algorithm,
  then $u$ cannot be deleted in the next $\sqrt{n}$ steps, since its degree
  falls by at most $1$ per step.
\end{observation}

Hence we have the following algorithm:
\begin{enumerate}
  \item Compute the set $S = \Set{u \in V \given \deg(u) > k + \sqrt{n}}$ of
    \emph{safe} vertices, and set it aside.
  \item Perform $\sqrt{n}$ iterations of the peeling algorithm, while updating
    the degree of a neighbor $u$ of a deleted vertex only if $u \notin S$ and $\deg(u) > k$.
  \item If the graph is non-empty, go to step 1.
\end{enumerate}

Let $N_A(u)$ and $\deg_A(u)$ denote the neighborhood and degree of $u$ within
$(V, E_A)$ (and similarly $N_B(u)$ and $\deg_B(u)$ for $E_B$).
We provide an implementation of the algorithm above in a communication protocol;
the sets $D$ and $L$ hold vertices that are ready to \emph{delete}, and
\emph{low} degree respectively.

\begin{Algorithm}\label{alg:nrootn}
  The $\Ot{n\sqrt{n}}$ communication protocol for $\degen_{n, k}$. \medskip

  \textbf{Input:} Alice and Bob get disjoint edge sets $E_A$ and $E_B$
    respectively, and both players get an integer $k$.\\
  \textbf{Output:} \true{} if $G = E_A \cup E_B$ has degeneracy $\leq k$, \fail{} otherwise.
  \begin{enumerate}
    \item Alice and Bob send $\deg_A(u)$ and $\deg_B(u)$ respectively for every
      vertex $u \in V(G)$, so they can then compute:
      \begin{itemize}
        \item $D \gets \Set*{ u \in G \given \deg(u) \leq k }$.
        \item $L \gets \Set*{ u \in G \given k + 1 \leq \deg(u) \leq k + \sqrt{n} }$.
        \item $S \gets V(G) \setminus (D \cup L)$.
      \end{itemize}
    \item For $\sqrt{n}$ iterations:
      \begin{enumerate}
        \item If $D$ is empty, output \fail.
        \item Let $v$ be any vertex in $D$. To remove $v$ from $G$:
        \item\label{step:Lnbrs} Alice and Bob send $N_A(v) \cap L$ and $N_B(v) \cap L$ respectively.
        \item\label{step:Lupdt} Hence they can update $\deg(u)$ for all $u \in L$.
        \item\label{step:Dmove} If $\deg(u) < k + 1$ for any $u \in L$,
          Alice and Bob move it to $D$.
      \end{enumerate}

    \item If $G$ is non-empty, go to Step 1.
    \item Output \true.
  \end{enumerate}
\end{Algorithm}
Note that the proofs (a $k$-ordering of the vertices in the \true{} case,
and a $(k + 1)$-core in the \fail{} case) fall out of the protocol exactly as before.
The correctness of this protocol follows from \Cref{obs:sqrt-local}; we omit
the proof here since it is subsumed by the next section.
To see that it only uses $\Ot{n \sqrt{n}}$ bits of communication, we observe
that:
\begin{itemize}
  \item Sending the degree of all vertices costs $O( n \log n )$ bits, and
    we do this at most $\sqrt{n}$ times.
  \item Once a vertex is in $D$, no communication is spent for updating its own degree.
  \item Each vertex in $L$ features at most $\sqrt{n}$ times in the neighborhood
    of a deleted vertex before it is moved to $D$, costing $O( \sqrt{n} \log n )$
    bits per vertex.
\end{itemize}
Hence the total communication cost of the protocol is $O( n\sqrt{n} \log n)$.

\subsection{An \texorpdfstring{$\Ot{n}$}{Õ(n)} Communication Protocol}

There are two changes from the previous protocol:
\begin{enumerate}
  \item 
    Instead of a single safe set of vertices, we have $O( \log n )$ safe sets
    corresponding to exponentially increasing gaps between the degree and $k$.
    More precisely, we define the set $S_i$ for $i \in \brac{\log n}$ such that
    whenever Alice and Bob compute $\deg(u)$ for some vertex $u$, they place it
    in $S_i$ if $k + 2^{i - 1} \leq \deg(u) < k + 2^i$.
  \item
    Instead of counting a global number of iterations until the safe vertices
    are updated, we track each vertex individually.
    Concretely, suppose $\deg(u) \geq k + \ell$.
    Then $u$ cannot be deleted until $\deg_A(u)$ or $\deg_B(u)$ falls by
    at least $\ell / 2$.
\end{enumerate}
The upshot of the two changes above is that for $u \in S_i$ Alice and Bob do not
have to update $\deg(u)$ until $\deg_A(u)$ or $\deg_B(u)$ falls by at least
$2^{i - 2}$.
We can now state the algorithm.
As before, the set $D$ holds all vertices which are ready for deletion.

\begin{Algorithm}\label{alg:n}
  The $\Ot{n}$ communication protocol for $\degen_{n, k}$. \medskip

  \textbf{Input:} Alice and Bob get disjoint edge sets $E_A$ and $E_B$
  respectively, and both players get an integer $k$.\\
  \textbf{Output:} \true{} if $G = E_A \cup E_B$ has degeneracy $\leq k$, \fail{} otherwise.
  \begin{enumerate}
    \item Alice and Bob send $\deg_A(u)$ and $\deg_B(u)$ respectively for every
      vertex $u \in V(G)$, so they can then compute:
      \begin{itemize}
        \item $D \gets \Set*{ u \in G \given \deg(u) \leq k }$.
        \item $S_i \gets \Set*{ u \in G \given k + 2^{i - 1} \leq \deg(u) < k + 2^i }$,
          for each $i \in \brac{\log n}$.
      \end{itemize}
    \item While {$G$ is non-empty:}
      \begin{enumerate}

        \item If $D$ is empty, output \fail.
        \item Let $v$ be any vertex in $D$. To remove $v$ from $G$:
        \item For each $i \in \brac{ \log n }$, and each $u \in S_i$,
          if $\deg_A(u)$ or $\deg_B(u)$ has fallen by $\geq 2^{i - 2}$ since it
          was last communicated, Alice sends $\deg_A(u)$ and Bob sends $\deg_B(u)$,
          hence allowing them to update $\deg(u)$.
        \item For all vertices updated in the previous step, Alice and Bob place
          them in the correct $S_i$ (or $D$).
      \end{enumerate}
    \item Output \true.
  \end{enumerate}
  %  \caption{The Peeling Algorithm for $\degen(G, k)$}\label{alg:peeling}
\end{Algorithm}

To wrap up, we prove that the algorithm above is correct, and uses $\Ot{n}$
space.

\begin{proof}[Proof of \Cref{lem:degen-decision}]
To prove that this protocol is correct, we argue that:
\begin{itemize}
  \item All vertices have degree $\leq k$ when they are removed, which
    implies that if the protocol outputs \true{}, $\kappa(G) \leq k$.
    This is immediate because we only move vertices to $D$ when we know their
    current degree exactly, and it is $\leq k$.
  \item Any vertex $u$ with degree $\leq k$ is placed in $D$, which implies
    that if $\kappa(G) \leq k$, then the protocol outputs \true{}.
    Suppose towards a contradiction that $u$ was placed in the set $S_i$ when
    its degree was updated for the final time.
    This means that $\deg(u)$ was at least $k + 2^{i - 1}$.
    But now for $\deg(u)$ to fall below $k + 1$, at least one of $\deg_A(u)$
    or $\deg_B(u)$ must fall by $\geq 2^{i - 2}$ (otherwise $\deg(u)$ falls by
    less than $2^{i - 1}$).
    Hence $\deg(u)$ was updated again following the final update, and we have
    a contradiction.
\end{itemize}

Next, we argue that this protocol uses only $\Ot{n}$ communication.
\begin{itemize}
  \item As before, the initial setup uses only $O(\log n)$ bits per vertex.
  \item While in $S_i$, a vertex can have its degree updated at most twice
    before it falls below $k + 2^{i - 1}$, and is hence dropped from
    $S_i$.
    Since there are only $\log n$ sets $S_i$, and a vertex can never gain
    degree, we only update its degree $O(\log n)$ times before it falls into $D$,
    spending $O( \log^2 n )$ bits overall.
\end{itemize}
Hence, in total, the protocol uses only $O(n \log^2 n)$ bits of communication.
\end{proof}

\newcommand{\etalchar}[1]{$^{#1}$}
\begin{thebibliography}{BvdBE{\etalchar{+}}22}

\bibitem[A17]{Assadi17sc}
Sepehr~Assadi A.
\newblock Tight space-approximation tradeoff for the multi-pass streaming set cover problem.
\newblock In {\em Proceedings of the 36th {ACM} {SIGMOD-SIGACT-SIGAI} Symposium on Principles of Database Systems, {PODS} 2017, Chicago, IL, USA, May 14-19, 2017}, pages 321--335, 2017.

\bibitem[A22]{Assadi22}
Sepehr~Assadi A.
\newblock A two-pass (conditional) lower bound for semi-streaming maximum matching.
\newblock In Joseph~(Seffi) Naor and Niv Buchbinder, editors, {\em Proceedings of the 2022 {ACM-SIAM} Symposium on Discrete Algorithms, {SODA} 2022, Virtual Conference / Alexandria, VA, USA, January 9 - 12, 2022}, pages 708--742. {SIAM}, 2022.

\bibitem[AA20]{AlonA20}
Noga Alon and Sepehr Assadi.
\newblock Palette sparsification beyond ({\(\Delta\)}+1) vertex coloring.
\newblock In Jaroslaw Byrka and Raghu Meka, editors, {\em Approximation, Randomization, and Combinatorial Optimization. Algorithms and Techniques, {APPROX/RANDOM} 2020, August 17-19, 2020, Virtual Conference}, volume 176 of {\em LIPIcs}, pages 6:1--6:22. Schloss Dagstuhl - Leibniz-Zentrum f{\"{u}}r Informatik, 2020.

\bibitem[ACK19]{AssadiCK19}
Sepehr Assadi, Yu~Chen, and Sanjeev Khanna.
\newblock Polynomial pass lower bounds for graph streaming algorithms.
\newblock In {\em {STOC}}, pages 265--276. {ACM}, 2019.

\bibitem[ACKP21]{AbboudCKP21}
Amir Abboud, Keren Censor{-}Hillel, Seri Khoury, and Ami Paz.
\newblock Smaller cuts, higher lower bounds.
\newblock {\em {ACM} Trans. Algorithms}, 17(4):30:1--30:40, 2021.

\bibitem[AD21]{AssadiD21}
Sepehr Assadi and Aditi Dudeja.
\newblock A simple semi-streaming algorithm for global minimum cuts.
\newblock In {\em {SOSA}}, pages 172--180. {SIAM}, 2021.

\bibitem[ADBV05]{Alvarez-HamelinDBV05}
J.~Ignacio Alvarez{-}Hamelin, Luca Dall'Asta, Alain Barrat, and Alessandro Vespignani.
\newblock Large scale networks fingerprinting and visualization using the k-core decomposition.
\newblock In {\em Advances in Neural Information Processing Systems 18 [Neural Information Processing Systems, {NIPS} 2005, December 5-8, 2005, Vancouver, British Columbia, Canada]}, pages 41--50, 2005.

\bibitem[AJJ{\etalchar{+}}22]{AssadiJJST22}
Sepehr Assadi, Arun Jambulapati, Yujia Jin, Aaron Sidford, and Kevin Tian.
\newblock Semi-streaming bipartite matching in fewer passes and optimal space.
\newblock In {\em {SODA}}, pages 627--669. {SIAM}, 2022.

\bibitem[AKSY20]{AssadiKSY20}
Sepehr Assadi, Gillat Kol, Raghuvansh Saxena, and Huacheng Yu.
\newblock Multi-pass graph streaming lower bounds for cycle counting, max-cut, matching size, and other problems.
\newblock In {\em 61st Annual {IEEE} Symposium on Foundations of Computer Science, {FOCS} (to appear)}, 2020.

\bibitem[AKZ22]{AssadiKZ22}
Sepehr Assadi, Gillat Kol, and Zhijun Zhang.
\newblock Rounds vs communication tradeoffs for maximal independent sets.
\newblock In {\em 63rd {IEEE} Annual Symposium on Foundations of Computer Science, {FOCS} 2022, Denver, CO, USA, October 31 - November 3, 2022}, pages 1193--1204. {IEEE}, 2022.

\bibitem[AN21]{AssadiN21}
Sepehr Assadi and Vishvajeet N.
\newblock Graph streaming lower bounds for parameter estimation and property testing via a streaming {XOR} lemma.
\newblock In Samir Khuller and Virginia~Vassilevska Williams, editors, {\em {STOC} '21: 53rd Annual {ACM} {SIGACT} Symposium on Theory of Computing, Virtual Event, Italy, June 21-25, 2021}, pages 612--625. {ACM}, 2021.

\bibitem[AR20]{AssadiR20}
Sepehr Assadi and Ran Raz.
\newblock Near-quadratic lower bounds for two-pass graph streaming algorithms.
\newblock In {\em {FOCS}}, pages 342--353. {IEEE}, 2020.

\bibitem[AYZ97]{AlonYZ97}
Noga Alon, Raphael Yuster, and Uri Zwick.
\newblock Finding and counting given length cycles.
\newblock {\em Algorithmica}, 17(3):209--223, 1997.

\bibitem[BBCR10]{BarakBCR10}
Boaz Barak, Mark Braverman, Xi~Chen, and Anup Rao.
\newblock How to compress interactive communication.
\newblock In {\em Proceedings of the 42nd {ACM} Symposium on Theory of Computing, {STOC} 2010, 5-8 June 2010}, pages 67--76, 2010.

\bibitem[BCD{\etalchar{+}}19]{BachrachCDELP19}
Nir Bachrach, Keren Censor{-}Hillel, Michal Dory, Yuval Efron, Dean Leitersdorf, and Ami Paz.
\newblock Hardness of distributed optimization.
\newblock In Peter Robinson and Faith Ellen, editors, {\em Proceedings of the 2019 {ACM} Symposium on Principles of Distributed Computing, {PODC} 2019, Toronto, ON, Canada, July 29 - August 2, 2019}, pages 238--247. {ACM}, 2019.

\bibitem[BCG20]{BeraCG19}
Suman~K. Bera, Amit Chakrabarti, and Prantar Ghosh.
\newblock Graph coloring via degeneracy in streaming and other space-conscious models.
\newblock In {\em 47th International Colloquium on Automata, Languages, and Programming, {ICALP} 2020, July 8-11, 2020, Saarbr{\"{u}}cken, Germany (Virtual Conference)}, pages 11:1--11:21, 2020.

\bibitem[BEO{\etalchar{+}}13]{BravermanEOPV13}
Mark Braverman, Faith Ellen, Rotem Oshman, Toniann Pitassi, and Vinod Vaikuntanathan.
\newblock A tight bound for set disjointness in the message-passing model.
\newblock In {\em 54th Annual {IEEE} Symposium on Foundations of Computer Science, {FOCS} 2013, 26-29 October, 2013, Berkeley, CA, {USA}}, pages 668--677, 2013.

\bibitem[BGKV14]{BonchiGKV14}
Francesco Bonchi, Francesco Gullo, Andreas Kaltenbrunner, and Yana Volkovich.
\newblock Core decomposition of uncertain graphs.
\newblock In Sofus~A. Macskassy, Claudia Perlich, Jure Leskovec, Wei Wang, and Rayid Ghani, editors, {\em The 20th {ACM} {SIGKDD} International Conference on Knowledge Discovery and Data Mining, {KDD} '14, New York, NY, {USA} - August 24 - 27, 2014}, pages 1316--1325. {ACM}, 2014.

\bibitem[BHNT15]{BhattacharyaHNT15}
Sayan Bhattacharya, Monika Henzinger, Danupon Nanongkai, and Charalampos~E. Tsourakakis.
\newblock Space- and time-efficient algorithm for maintaining dense subgraphs on one-pass dynamic streams.
\newblock In {\em Proceedings of the Forty-Seventh Annual {ACM} on Symposium on Theory of Computing, {STOC} 2015, Portland, OR, USA, June 14-17, 2015}, pages 173--182, 2015.

\bibitem[BJKS02]{Bar-YossefJKS02}
Ziv Bar{-}Yossef, T.~S. Jayram, Ravi Kumar, and D.~Sivakumar.
\newblock An information statistics approach to data stream and communication complexity.
\newblock In {\em 43rd Symposium on Foundations of Computer Science {(FOCS} 2002), 16-19 November 2002, Proceedings}, pages 209--218, 2002.

\bibitem[BKS02]{Bar-YossefKS02}
Ziv Bar{-}Yossef, Ravi Kumar, and D.~Sivakumar.
\newblock Reductions in streaming algorithms, with an application to counting triangles in graphs.
\newblock In {\em Proceedings of the Thirteenth Annual {ACM-SIAM} Symposium on Discrete Algorithms, January 6-8, 2002, San Francisco, CA, {USA.}}, pages 623--632, 2002.

\bibitem[BKV12]{BahmaniKV12}
Bahman Bahmani, Ravi Kumar, and Sergei Vassilvitskii.
\newblock Densest subgraph in streaming and mapreduce.
\newblock {\em {PVLDB}}, 5(5):454--465, 2012.

\bibitem[BR11]{BravermanR11}
Mark Braverman and Anup Rao.
\newblock Information equals amortized communication.
\newblock In {\em {IEEE} 52nd Annual Symposium on Foundations of Computer Science, {FOCS} 2011, October 22-25, 2011}, pages 748--757, 2011.

\bibitem[BvdBE{\etalchar{+}}22]{blikstadBEMN22}
Joakim Blikstad, Jan van~den Brand, Yuval Efron, Sagnik Mukhopadhyay, and Danupon Nanongkai.
\newblock Nearly optimal communication and query complexity of bipartite matching.
\newblock In {\em {FOCS}}, pages 1174--1185. {IEEE}, 2022.

\bibitem[CFHT20]{ChangFHT20}
Yi{-}Jun Chang, Martin Farach{-}Colton, Tsan{-}sheng Hsu, and Meng{-}Tsung Tsai.
\newblock Streaming complexity of spanning tree computation.
\newblock In {\em 37th International Symposium on Theoretical Aspects of Computer Science, {STACS} 2020, March 10-13, 2020, Montpellier, France}, pages 34:1--34:19, 2020.

\bibitem[CGMV20]{ChakrabartiGMV20}
Amit Chakrabarti, Prantar Ghosh, Andrew McGregor, and Sofya Vorotnikova.
\newblock Vertex ordering problems in directed graph streams.
\newblock In {\em Proceedings of the 2020 {ACM-SIAM} Symposium on Discrete Algorithms, {SODA} 2020, Salt Lake City, UT, USA, January 5-8, 2020}, pages 1786--1802, 2020.

\bibitem[Cha00]{Charikar00}
Moses Charikar.
\newblock Greedy approximation algorithms for finding dense components in a graph.
\newblock In Klaus Jansen and Samir Khuller, editors, {\em Approximation Algorithms for Combinatorial Optimization, Third International Workshop, {APPROX} 2000, Saarbr{\"{u}}cken, Germany, September 5-8, 2000, Proceedings}, volume 1913 of {\em Lecture Notes in Computer Science}, pages 84--95. Springer, 2000.

\bibitem[CKP{\etalchar{+}}21a]{ChenKPSSY21}
Lijie Chen, Gillat Kol, Dmitry Paramonov, Raghuvansh~R. Saxena, Zhao Song, and Huacheng Yu.
\newblock Almost optimal super-constant-pass streaming lower bounds for reachability.
\newblock In Samir Khuller and Virginia~Vassilevska Williams, editors, {\em {STOC} '21: 53rd Annual {ACM} {SIGACT} Symposium on Theory of Computing, Virtual Event, Italy, June 21-25, 2021}, pages 570--583. {ACM}, 2021.

\bibitem[CKP{\etalchar{+}}21b]{ChenKPSSY21b}
Lijie Chen, Gillat Kol, Dmitry Paramonov, Raghuvansh~R. Saxena, Zhao Song, and Huacheng Yu.
\newblock Near-optimal two-pass streaming algorithm for sampling random walks over directed graphs.
\newblock In Nikhil Bansal, Emanuela Merelli, and James Worrell, editors, {\em 48th International Colloquium on Automata, Languages, and Programming, {ICALP} 2021, July 12-16, 2021, Glasgow, Scotland (Virtual Conference)}, volume 198 of {\em LIPIcs}, pages 52:1--52:19. Schloss Dagstuhl - Leibniz-Zentrum f{\"{u}}r Informatik, 2021.

\bibitem[CKP{\etalchar{+}}23]{ChenKPSSY23}
Lijie Chen, Gillat Kol, Dmitry Paramonov, Raghuvansh~R. Saxena, Zhao Song, and Huacheng Yu.
\newblock Towards multi-pass streaming lower bounds for optimal approximation of max-cut.
\newblock In Nikhil Bansal and Viswanath Nagarajan, editors, {\em Proceedings of the 2023 {ACM-SIAM} Symposium on Discrete Algorithms, {SODA} 2023, Florence, Italy, January 22-25, 2023}, pages 878--924. {SIAM}, 2023.

\bibitem[CSWY01]{ChakrabartiSWY01}
Amit Chakrabarti, Yaoyun Shi, Anthony Wirth, and Andrew~Chi{-}Chih Yao.
\newblock Informational complexity and the direct sum problem for simultaneous message complexity.
\newblock In {\em 42nd Annual Symposium on Foundations of Computer Science, {FOCS} 2001, 14-17 October 2001}, pages 270--278, 2001.

\bibitem[CT06]{CoverT06}
Thomas~M. Cover and Joy~A. Thomas.
\newblock {\em Elements of information theory {(2.} ed.)}.
\newblock Wiley, 2006.

\bibitem[CZL{\etalchar{+}}20]{ChuZ00ZXZ20}
Deming Chu, Fan Zhang, Xuemin Lin, Wenjie Zhang, Ying Zhang, Yinglong Xia, and Chenyi Zhang.
\newblock Finding the best k in core decomposition: {A} time and space optimal solution.
\newblock In {\em 36th {IEEE} International Conference on Data Engineering, {ICDE} 2020, Dallas, TX, USA, April 20-24, 2020}, pages 685--696. {IEEE}, 2020.

\bibitem[DBS17]{DhulipalaBS17}
Laxman Dhulipala, Guy~E. Blelloch, and Julian Shun.
\newblock Julienne: {A} framework for parallel graph algorithms using work-efficient bucketing.
\newblock In Christian Scheideler and Mohammad~Taghi Hajiaghayi, editors, {\em Proceedings of the 29th {ACM} Symposium on Parallelism in Algorithms and Architectures, {SPAA} 2017, Washington DC, USA, July 24-26, 2017}, pages 293--304. {ACM}, 2017.

\bibitem[DBS18]{DhulipalaBS18}
Laxman Dhulipala, Guy~E. Blelloch, and Julian Shun.
\newblock Theoretically efficient parallel graph algorithms can be fast and scalable.
\newblock In Christian Scheideler and Jeremy~T. Fineman, editors, {\em Proceedings of the 30th on Symposium on Parallelism in Algorithms and Architectures, {SPAA} 2018, Vienna, Austria, July 16-18, 2018}, pages 393--404. {ACM}, 2018.

\bibitem[DNO19]{DobzinskiNO19}
Shahar Dobzinski, Noam Nisan, and Sigal Oren.
\newblock Economic efficiency requires interaction.
\newblock {\em Games Econ. Behav.}, 118:589--608, 2019.

\bibitem[EH66]{erdHos1966chromatic}
Paul Erd{\H{o}}s and Andr{\'a}s Hajnal.
\newblock On chromatic number of graphs and set-systems.
\newblock {\em Acta Math. Acad. Sci. Hungar}, 17(61-99):1, 1966.

\bibitem[ELM18]{EsfandiariLM18}
Hossein Esfandiari, Silvio Lattanzi, and Vahab~S. Mirrokni.
\newblock Parallel and streaming algorithms for k-core decomposition.
\newblock In Jennifer~G. Dy and Andreas Krause, editors, {\em Proceedings of the 35th International Conference on Machine Learning, {ICML} 2018, Stockholmsm{\"{a}}ssan, Stockholm, Sweden, July 10-15, 2018}, volume~80 of {\em Proceedings of Machine Learning Research}, pages 1396--1405. {PMLR}, 2018.

\bibitem[FHW12]{FrischknechtHW12}
Silvio Frischknecht, Stephan Holzer, and Roger Wattenhofer.
\newblock Networks cannot compute their diameter in sublinear time.
\newblock In Yuval Rabani, editor, {\em Proceedings of the Twenty-Third Annual {ACM-SIAM} Symposium on Discrete Algorithms, {SODA} 2012, Kyoto, Japan, January 17-19, 2012}, pages 1150--1162. {SIAM}, 2012.

\bibitem[FKM{\etalchar{+}}05]{FeigenbaumKMSZ05}
Joan Feigenbaum, Sampath Kannan, Andrew McGregor, Siddharth Suri, and Jian Zhang.
\newblock On graph problems in a semi-streaming model.
\newblock {\em Theor. Comput. Sci.}, 348(2-3):207--216, 2005.

\bibitem[FT14]{Farach-ColtonT14}
Martin Farach{-}Colton and Meng{-}Tsung Tsai.
\newblock Computing the degeneracy of large graphs.
\newblock In Alberto Pardo and Alfredo Viola, editors, {\em {LATIN} 2014: Theoretical Informatics - 11th Latin American Symposium, Montevideo, Uruguay, March 31 - April 4, 2014. Proceedings}, volume 8392 of {\em Lecture Notes in Computer Science}, pages 250--260. Springer, 2014.

\bibitem[FT16]{Farach-ColtonT16}
Martin Farach{-}Colton and Meng{-}Tsung Tsai.
\newblock Tight approximations of degeneracy in large graphs.
\newblock In {\em {LATIN} 2016: Theoretical Informatics - 12th Latin American Symposium, Ensenada, Mexico, April 11-15, 2016, Proceedings}, pages 429--440, 2016.

\bibitem[GBGL20]{GalimbertiBGL20}
Edoardo Galimberti, Francesco Bonchi, Francesco Gullo, and Tommaso Lanciano.
\newblock Core decomposition in multilayer networks: Theory, algorithms, and applications.
\newblock {\em {ACM} Trans. Knowl. Discov. Data}, 14(1):11:1--11:40, 2020.

\bibitem[GLM19]{GhaffariLM19}
Mohsen Ghaffari, Silvio Lattanzi, and Slobodan Mitrovic.
\newblock Improved parallel algorithms for density-based network clustering.
\newblock In Kamalika Chaudhuri and Ruslan Salakhutdinov, editors, {\em Proceedings of the 36th International Conference on Machine Learning, {ICML} 2019, 9-15 June 2019, Long Beach, California, {USA}}, volume~97 of {\em Proceedings of Machine Learning Research}, pages 2201--2210. {PMLR}, 2019.

\bibitem[GO13]{GuruswamiO13}
Venkatesan Guruswami and Krzysztof Onak.
\newblock Superlinear lower bounds for multipass graph processing.
\newblock In {\em Proceedings of the 28th Conference on Computational Complexity, {CCC} 2013, K.lo Alto, California, USA, 5-7 June, 2013}, pages 287--298, 2013.

\bibitem[HSSW12]{HalldorssonSSW12}
Magn{\'{u}}s~M. Halld{\'{o}}rsson, Xiaoming Sun, Mario Szegedy, and Chengu Wang.
\newblock Streaming and communication complexity of clique approximation.
\newblock In {\em Automata, Languages, and Programming - 39th International Colloquium, {ICALP} 2012, Warwick, UK, July 9-13, 2012, Proceedings, Part {I}}, pages 449--460, 2012.

\bibitem[IKL{\etalchar{+}}12]{IvanyosKLSW12}
G{\'{a}}bor Ivanyos, Hartmut Klauck, Troy Lee, Miklos Santha, and Ronald de~Wolf.
\newblock New bounds on the classical and quantum communication complexity of some graph properties.
\newblock In {\em {FSTTCS}}, volume~18 of {\em LIPIcs}, pages 148--159. Schloss Dagstuhl - Leibniz-Zentrum f{\"{u}}r Informatik, 2012.

\bibitem[JKS03]{JayramKS03}
T.~S. Jayram, Ravi Kumar, and D.~Sivakumar.
\newblock Two applications of information complexity.
\newblock In {\em {STOC}}, pages 673--682. {ACM}, 2003.

\bibitem[KBVT15]{KhaouidBST15}
Wissam Khaouid, Marina Barsky, S.~Venkatesh, and Alex Thomo.
\newblock K-core decomposition of large networks on a single {PC}.
\newblock {\em Proc. {VLDB} Endow.}, 9(1):13--23, 2015.

\bibitem[Kla00]{klauck2000quantum}
Hartmut Klauck.
\newblock On quantum and probabilistic communication: Las vegas and one-way protocols.
\newblock In {\em 32nd Annual ACM Symposium on Theory of Computing ({STOC})}, pages 644--651, 2000.

\bibitem[KN97]{KushilevitzN97}
Eyal Kushilevitz and Noam Nisan.
\newblock {\em Communication complexity}.
\newblock Cambridge University Press, 1997.

\bibitem[KPP16]{KopelowitzPP16}
Tsvi Kopelowitz, Seth Pettie, and Ely Porat.
\newblock Higher lower bounds from the 3sum conjecture.
\newblock In Robert Krauthgamer, editor, {\em Proceedings of the Twenty-Seventh Annual {ACM-SIAM} Symposium on Discrete Algorithms, {SODA} 2016, Arlington, VA, USA, January 10-12, 2016}, pages 1272--1287. {SIAM}, 2016.

\bibitem[KPSY23]{KolPSY23}
Gillat Kol, Dmitry Paramonov, Raghuvansh~R. Saxena, and Huacheng Yu.
\newblock Characterizing the multi-pass streaming complexity for solving boolean csps exactly.
\newblock In Yael~Tauman Kalai, editor, {\em 14th Innovations in Theoretical Computer Science Conference, {ITCS} 2023, January 10-13, 2023, MIT, Cambridge, Massachusetts, {USA}}, volume 251 of {\em LIPIcs}, pages 80:1--80:15. Schloss Dagstuhl - Leibniz-Zentrum f{\"{u}}r Informatik, 2023.

\bibitem[LJS19]{LiuJS19}
Yang~P. Liu, Arun Jambulapati, and Aaron Sidford.
\newblock Parallel reachability in almost linear work and square root depth.
\newblock In David Zuckerman, editor, {\em 60th {IEEE} Annual Symposium on Foundations of Computer Science, {FOCS} 2019, Baltimore, Maryland, USA, November 9-12, 2019}, pages 1664--1686. {IEEE} Computer Society, 2019.

\bibitem[LSY{\etalchar{+}}22]{LiuSYDS22}
Quanquan~C. Liu, Jessica Shi, Shangdi Yu, Laxman Dhulipala, and Julian Shun.
\newblock Parallel batch-dynamic algorithms for k-core decomposition and related graph problems.
\newblock In Kunal Agrawal and I{-}Ting~Angelina Lee, editors, {\em {SPAA} '22: 34th {ACM} Symposium on Parallelism in Algorithms and Architectures, Philadelphia, PA, USA, July 11 - 14, 2022}, pages 191--204. {ACM}, 2022.

\bibitem[LWSX18]{LiWSX18}
Chao Li, Li~Wang, Shiwen Sun, and Chengyi Xia.
\newblock Identification of influential spreaders based on classified neighbors in real-world complex networks.
\newblock {\em Appl. Math. Comput.}, 320:512--523, 2018.

\bibitem[LZZ{\etalchar{+}}19]{LiZZQZL19}
Conggai Li, Fan Zhang, Ying Zhang, Lu~Qin, Wenjie Zhang, and Xuemin Lin.
\newblock Efficient progressive minimum k-core search.
\newblock {\em Proc. {VLDB} Endow.}, 13(3):362--375, 2019.

\bibitem[MB83]{MatulaB83}
David~W. Matula and Leland~L. Beck.
\newblock Smallest-last ordering and clustering and graph coloring algorithms.
\newblock {\em J. {ACM}}, 30(3):417--427, 1983.

\bibitem[McG14]{McGregor14}
Andrew McGregor.
\newblock Graph stream algorithms: a survey.
\newblock {\em {SIGMOD} Rec.}, 43(1):9--20, 2014.

\bibitem[MN20]{MukhopadhyayN20}
Sagnik Mukhopadhyay and Danupon Nanongkai.
\newblock Weighted min-cut: sequential, cut-query, and streaming algorithms.
\newblock In {\em {STOC}}, pages 496--509. {ACM}, 2020.

\bibitem[MTVV15]{McGregorTVV15}
Andrew McGregor, David Tench, Sofya Vorotnikova, and Hoa~T. Vu.
\newblock Densest subgraph in dynamic graph streams.
\newblock In {\em Mathematical Foundations of Computer Science 2015 - 40th International Symposium, {MFCS} 2015, Milan, Italy, August 24-28, 2015, Proceedings, Part {II}}, pages 472--482, 2015.

\bibitem[NW91]{NisanW91}
Noam Nisan and Avi Wigderson.
\newblock Rounds in communication complexity revisited.
\newblock In {\em Proceedings of the 23rd Annual {ACM} Symposium on Theory of Computing, May 5-8, 1991, New Orleans, Louisiana, {USA}}, pages 419--429, 1991.

\bibitem[NW93]{NisanW93}
Noam Nisan and Avi Wigderson.
\newblock Rounds in communication complexity revisited.
\newblock {\em {SIAM} J. Comput.}, 22(1):211--219, 1993.

\bibitem[PRV99]{PonzioRV99}
Stephen Ponzio, Jaikumar Radhakrishnan, and Srinivasan Venkatesh.
\newblock The communication complexity of pointer chasing: Applications of entropy and sampling.
\newblock In {\em Proceedings of the Thirty-First Annual {ACM} Symposium on Theory of Computing, May 1-4, 1999, Atlanta, Georgia, {USA}}, pages 602--611, 1999.

\bibitem[PS84]{PapadimitriouS84}
Christos~H. Papadimitriou and Michael Sipser.
\newblock Communication complexity.
\newblock {\em J. Comput. Syst. Sci.}, 28(2):260--269, 1984.

\bibitem[Raz92]{Razborov92}
Alexander~A. Razborov.
\newblock On the distributional complexity of disjointness.
\newblock {\em Theor. Comput. Sci.}, 106(2):385--390, 1992.

\bibitem[RSW18]{RubinsteinSW18}
Aviad Rubinstein, Tselil Schramm, and S.~Matthew Weinberg.
\newblock Computing exact minimum cuts without knowing the graph.
\newblock In {\em {ITCS}}, volume~94 of {\em LIPIcs}, pages 39:1--39:16. Schloss Dagstuhl - Leibniz-Zentrum f{\"{u}}r Informatik, 2018.

\bibitem[RW13]{RodittyW13}
Liam Roditty and Virginia~Vassilevska Williams.
\newblock Fast approximation algorithms for the diameter and radius of sparse graphs.
\newblock In Dan Boneh, Tim Roughgarden, and Joan Feigenbaum, editors, {\em Symposium on Theory of Computing Conference, STOC'13, Palo Alto, CA, USA, June 1-4, 2013}, pages 515--524. {ACM}, 2013.

\bibitem[SGJ{\etalchar{+}}13]{SariyuceGJWC13}
Ahmet~Erdem Sariy{\"{u}}ce, Bugra Gedik, Gabriela Jacques{-}Silva, Kun{-}Lung Wu, and {\"{U}}mit~V. {\c{C}}ataly{\"{u}}rek.
\newblock Streaming algorithms for k-core decomposition.
\newblock {\em Proc. {VLDB} Endow.}, 6(6):433--444, 2013.

\bibitem[Yao79]{Yao79}
Andrew~Chi{-}Chih Yao.
\newblock Some complexity questions related to distributive computing (preliminary report).
\newblock In {\em Proceedings of the 11h Annual {ACM} Symposium on Theory of Computing, April 30 - May 2, 1979, Atlanta, Georgia, {USA}}, pages 209--213, 1979.

\bibitem[Yeh20]{Yehudayoff20}
Amir Yehudayoff.
\newblock Pointer chasing via triangular discrimination.
\newblock {\em Comb. Probab. Comput.}, 29(4):485--494, 2020.

\end{thebibliography}
\end{document}
